\documentclass[11pt]{article}
\usepackage[T1]{fontenc}
\usepackage[utf8]{inputenc}
\usepackage{lmodern}
\usepackage[margin=1in]{geometry}
\usepackage{amsmath,amssymb,amsthm}
\usepackage{microtype}
\usepackage{hyperref}
\hypersetup{colorlinks=true,linkcolor=blue,citecolor=blue,urlcolor=blue}
\newtheorem{theorem}{Theorem}
\newtheorem{lemma}{Lemma}
\newcommand{\Tr}{\operatorname{Tr}}
\newcommand{\ket}[1]{\lvert #1\rangle}
\newcommand{\bra}[1]{\langle #1\rvert}
\newcommand{\C}{\mathbb C}
\title{The sharp overlap threshold for antidistinguishability of five pure states}
\author{Xiaojian Zeng\\
\small PBC School of Finance, Tsinghua University, 100190\\
\small \texttt{xiaojianzeng@foxmail.com}}
\date{}
\begin{document}
\maketitle
\begin{abstract}
Five pure states with every pairwise overlap modulus at most
$\sqrt{3/8}$ admit a zero-error exclusion measurement by a known universal
criterion. We prove that this guarantee is sharp. For every
$3/8<m<2/5$, we construct five normalized complex states of rank three,
each pair having squared overlap $m$, and exhibit a finite Hermitian
dual certificate with strictly positive trace. The certificate excludes
every zero-error positive operator-valued measurement and supplies a
positive lower bound on the minimum exclusion error for equal priors.
Embedding the construction in $\C^5$ and taking $m$ down to $3/8$
establishes that the infimum of the maximum overlap over non-excludable
five-state ensembles is exactly $\sqrt{3/8}$. The boundary is included in
the sufficient criterion, so the infimum is not attained by a
non-excludable ensemble.
\end{abstract}

\section{Zero-error exclusion and the sharpness question}
In quantum state exclusion, outcome $i$ announces that the prepared
state was not $\psi_i$. For normalized pure states
$\psi_1,\ldots,\psi_5\in\C^5$, a zero-error exclusion measurement is a
positive operator-valued measurement (POVM) satisfying
\begin{equation}
 M_i\succeq0,\qquad \sum_{i=1}^5 M_i=I_5,\qquad
 \bra{\psi_i}M_i\ket{\psi_i}=0\quad (1\leq i\leq5).
 \label{eq:povm}
\end{equation}
We call such a tuple antidistinguishable. Linear dependence among the
states and zero measurement effects are allowed. The problem is
single-copy exclusion with arbitrary PSD effects, without postselection
or restrictions to projective measurements.

Let
\begin{equation}
 \mu(\psi_1,\ldots,\psi_5)=
 \max_{i<j}|\langle\psi_i\mid\psi_j\rangle|,\qquad
 t_5=\inf_{\text{non-antidistinguishable tuples}}\mu .
 \label{eq:threshold}
\end{equation}
Johnston, Russo, and Sikora proved that $n$ pure states are
antidistinguishable whenever their maximum overlap is at most
$\sqrt{(n-2)/(2n-2)}$ \cite[Corollary 5.4]{JRS25}.
Their Section 6 establishes sharpness through four states and leaves
the higher-state question open. The subsequent classification framework
of Johnston, Lovitz, Russo, and Sikora retains an open entrywise
sharpness question, while proving sharpness of a distinct
Frobenius-norm condition \cite[Section 1.1 and Section 3]{JLRS26}.
Our contribution is the matching construction for five states.

\begin{theorem}\label{thm:main}
For the model in \eqref{eq:povm}, $t_5=\sqrt{3/8}$.
For each real $m\in(3/8,2/5)$, there is a rank-three tuple with
$|\langle\psi_i\mid\psi_j\rangle|^2=m$ for every $i\neq j$
which is not antidistinguishable. The explicit certificate below
also gives a uniform-prior exclusion error lower bound $s\eta/5>0$.
\end{theorem}

The family and the certificate are exact throughout this open
interval. The result resolves the five-state sector; it makes no
assertion about sharpness for six or more states, an exact minimum-error
formula, or experimental robustness.

\section{A rank-three equimodular family}
Fix $3/8<m<2/5$ and define
\begin{equation}
 p=\frac{1+\sqrt{(4m-1)/3}}2,\quad q=\frac mp,\quad
 h=\frac{\sqrt{m-(2q-1)^2}}2,\quad
 k=\frac{\sqrt{1-m}}2,\quad \omega=e^{2\pi i/3}.
 \label{eq:parameters}
\end{equation}
All square roots denote positive real roots. In an orthonormal basis
$e_1,e_2,e_3$ of $\C^3$, set
\begin{align}
 \ket{a_j}&=\sqrt p\,e_1+\sqrt{1-p}\,\omega^j e_3
       &&(j=0,1,2),\label{eq:orbit}\\
 \ket{b_\pm}&=\sqrt q\,e_1+
                \frac{h\pm ik}{\sqrt q}\,e_2 .
       \label{eq:pair}
\end{align}
These five vectors, with two additional zero coordinates, are the
desired states in $\C^5$.

\begin{lemma}\label{lem:states}
The vectors \eqref{eq:orbit}--\eqref{eq:pair} are normalized,
span $\C^3$, and have squared pairwise overlap $m$.
\end{lemma}
\begin{proof}
Equation \eqref{eq:parameters} gives
$m=1-3p+3p^2$ and $7/10<p<3/4$. Hence
$1/2<q<4/7$ and
$0<(2q-1)^2<1/49<3/8<m<1$. All defining roots and
denominators are legal. Moreover,
\begin{equation}
 h^2+k^2=q(1-q).
 \label{eq:norm}
\end{equation}
This proves normalization. For distinct orbit states the phase
difference has real part $-1/2$, giving
\begin{equation}
 |\langle a_i\mid a_j\rangle|^2
   =p^2+(1-p)^2-p(1-p)=m.
\end{equation}
Each orbit-to-pair squared overlap is $pq=m$. For the final pair,
\begin{equation}
 \langle b_+\mid b_-\rangle=q+\frac{(h-ik)^2}{q},
 \qquad
 \left|q+\frac{(h-ik)^2}{q}\right|^2=1-4k^2=m,
 \label{eq:pair-overlap}
\end{equation}
where the second equality uses \eqref{eq:norm}.
Two orbit states span $e_1,e_3$; either pair state has a nonzero
$e_2$ component. Thus the rank is exactly three.
\end{proof}

\section{A finite exclusion certificate}
For any states with projectors $\rho_i$, a Hermitian operator $Y$
satisfying $Y\preceq\rho_i$ for all $i$ obeys
\begin{equation}
 \sum_i\Tr(M_i\rho_i)
  =\Tr Y+\sum_i\Tr\!\left[M_i(\rho_i-Y)\right]\geq\Tr Y
 \label{eq:weak-duality}
\end{equation}
for every POVM. This is direct weak duality for the exclusion
optimization \cite[Eq.~(2)]{JRS25}. A positive trace rules out
\eqref{eq:povm}. No strong-duality assumption is needed.

For the parameters in \eqref{eq:parameters}, define
\begin{align}
 r&=\frac{2p-1}{1-p},&
 x&=\frac{rh}{1-q},&
 L&=\frac{1+x^2}{r},\nonumber\\
 R&=\frac{2q-1+2hx}{1-q},&
 D&=R-L,&
 V&=\frac{L+R}{2}.
 \label{eq:auxiliary}
\end{align}
The positive gap driving the certificate is
\begin{equation}
 D=\frac{p^3(8m-3)}
          {4(1-p)^2(2p-1)(3p-1)^2}>0.
 \label{eq:gap}
\end{equation}
For completeness, substitution of $x$ first reduces it to
\[
 D=\frac{2q-1}{1-q}-\frac1r+\frac{rh^2}{(1-q)^2}.
\]
Using $q=m/p$, $h^2=[m-(2q-1)^2]/4$, and
$1-q=(1-p)(3p-1)/p$, the numerator of the common-denominator
expression becomes $p^3(24p^2-24p+5)=p^3(8m-3)$,
which proves \eqref{eq:gap}. In particular $r,x,L,R,V$ are positive.

Now put
\begin{align}
 \eta&=\frac{rD}{4(1+r)},&
 A&=1+V+\eta,\nonumber\\
 C_1&=\frac{pV^2}{(1-p)^2},&
 C_2&=\frac{x^2}{1-q}+\frac{q+2hx}{(1-q)^2},\nonumber\\
 s&=\frac{\eta}{2(C_1+C_2+1)},&
 Y&=s\begin{pmatrix}A&x&0\\x&-1&0\\0&0&-V\end{pmatrix}.
 \label{eq:dual}
\end{align}
These are finite quantities at each point of the open interval, and
\begin{equation}
 \Tr Y=s(A-1-V)=s\eta>0.
 \label{eq:trace}
\end{equation}
The operator $Y$ acts on the physical three-dimensional state span.
It is distinct from a dual witness acting on the five state labels
of a Gram matrix.

\begin{lemma}\label{lem:dual}
For each of the five states in \eqref{eq:orbit}--\eqref{eq:pair},
$\rho_i-Y$ is positive definite on $\C^3$.
\end{lemma}
\begin{proof}
For an orbit state, the lower $e_2,e_3$ block of $\rho_i-Y$ is
$\operatorname{diag}(s,1-p+sV)\succ0$. Its Schur complement is
\begin{equation}
 s[F_1(s)-A],\qquad F_1(s)=\frac{pV}{1-p+sV}-x^2.
 \label{eq:schur1}
\end{equation}
The orbit phase cancels because $|\omega^j|=1$.
For either pair state, the lower block is
$\operatorname{diag}(1-q+s,sV)\succ0$, and the Schur complement is
\begin{equation}
 s[F_2(s)-A],\qquad
 F_2(s)=\frac{q+2hx-sx^2}{1-q+s}.
 \label{eq:schur2}
\end{equation}
The sign of $k$ drops out since $x$ is real and
$h^2+k^2=q(1-q)$.

At zero step, the exact slacks are
\begin{equation}
 F_1(0)-A=\frac{rD}{2}-\eta>\eta,\qquad
 F_2(0)-A=\frac{D}{2}-\eta>\eta.
 \label{eq:initial}
\end{equation}
Indeed, subtracting another $\eta$ from the two expressions gives
$r^2D/[2(1+r)]$ and $D/[2(1+r)]$, respectively.
Both are strictly positive without any ordering assumption on $r$.
The finite changes satisfy the exact identities and bounds
\begin{align}
 F_1(0)-F_1(s)
   &=\frac{spV^2}{(1-p)(1-p+sV)}\leq sC_1,
       \label{eq:loss1}\\
 F_2(0)-F_2(s)
   &=\frac{s[(q+2hx)/(1-q)+x^2]}{1-q+s}\leq sC_2.
       \label{eq:loss2}
\end{align}
The step in \eqref{eq:dual} gives $sC_i<\eta/2$ for
$i=1,2$. Equations \eqref{eq:initial}--\eqref{eq:loss2}
therefore imply $F_i(s)-A>\eta/2>0$.
Both Schur complements are strictly positive, proving the lemma.
The bounds hold at a finite step and contain no omitted expansion
remainder.
\end{proof}

\section{The threshold and its operational meaning}
\begin{proof}[Proof of Theorem \ref{thm:main}]
Lemma \ref{lem:dual} and \eqref{eq:weak-duality} imply
$\sum_i\Tr(M_i\rho_i)\geq s\eta>0$ for every POVM on $\C^3$.
After embedding, $Y\oplus0_2\preceq\rho_i\oplus0_2$ on $\C^5$,
so the same bound applies to all POVMs in the original model.
This also covers measurements with zero effects. With equal priors,
\begin{equation}
 P_{\mathrm{err}}^{\mathrm{opt}}
  :=\min_{\{M_i\}}\frac15\sum_i\Tr(M_i\rho_i)
  \ \geq\ \frac{s\eta}{5}>0 .
 \label{eq:error}
\end{equation}
Lemma \ref{lem:states} supplies maximum overlap $\sqrt m$.
For every $\varepsilon>0$, choose
$3/8<m<\min\{2/5,(\sqrt{3/8}+\varepsilon)^2\}$.
The resulting non-antidistinguishable tuple gives
$t_5\leq\sqrt{3/8}$.

Conversely, the published overlap criterion
\cite[Corollary 5.4]{JRS25} at $n=5$ guarantees
antidistinguishability when $\mu\leq\sqrt{3/8}$.
Equivalently, its sufficient Frobenius condition applies because
$\|G\|_F^2\leq5+20(3/8)=25/2$ for the Gram matrix
$G_{ij}=\langle\psi_i\mid\psi_j\rangle$.
It follows that $t_5\geq\sqrt{3/8}$, completing the proof.
\end{proof}

The threshold determines what can be guaranteed using overlap
magnitudes alone. Every five-state ensemble at or below it admits
perfect exclusion, while arbitrarily close above it there are legal
ensembles for which every exclusion measurement makes errors.
The construction uses complex phases and a three-dimensional span;
the operational obstruction is the positive trace of the feasible
physical dual, rather than failure of a sufficient overlap test.

The family approaches an excludable boundary. The certificate's
positive trace is needed only at $m>3/8$, and \eqref{eq:error}
is a lower bound rather than the exact optimum. We do not obtain
a non-excludable ensemble at equality, a universal lower bound on
error for other ensembles, or a noise tolerance.

\section{Reproducibility and scope}
The accompanying exact SymPy checker verifies the rational identity
\eqref{eq:gap}, normalization and overlap relations, initial Schur
slacks, and finite-step loss identities. It also constructs the five
projectors and the dual directly at three rational values of $p$,
checking all nonempty principal minors of each $\rho_i-Y$, all ten
overlaps, and the positive trace. The analytic sign and Schur argument
establishes the whole interval; those finite checks are supplementary.
The proof does not use numerical optimization, an SDP solver status,
or a perturbative approximation. All vectors, effects, and overlaps
are dimensionless ideal quantum-state quantities.

The literature comparison is restricted to the cited primary works
and a targeted search through 12 October 2026; no exhaustive priority
claim is intended. In particular, sharpness of a Frobenius bound does
not by itself establish the maximum-entry overlap threshold considered
here. The conclusions concern unrestricted single-copy pure-state
exclusion and do not extend to a different measurement convention
or a multiple-copy task.

\section*{Author contribution and AI disclosure}
OpenAI Codex was used in developing the construction and proof,
symbolic and numerical exploratory calculations, code generation,
literature searches, internal verification, and manuscript drafting.
The internal checks used AI-assisted implementations and do not
constitute external peer review. Human scientific review, confirmation
of authorship and affiliation, and approval of the final manuscript
remain pending for this local draft.

\section*{Code availability}
The accompanying source package contains an exact verification script,
its recorded output, and instructions identifying the required SymPy
version. No experimental data are used.

\bibliographystyle{plain}
\bibliography{references}
\end{document}
